\chapter{The SPI seam}
\label{ch:spi}

\chapterbrief
  {Define what ``an SPI bus'' means to your program, build a fake one, and test against it with no hardware attached.}
  {Interfaces, and dependency injection.}
  {A green test that runs with the reader unplugged.}

\section{The problem}

You are about to write a driver that pokes RC522 registers. Before you write a
line of it, ask an uncomfortable question:

\begin{quote}\itshape How will I know the driver is correct?\end{quote}

The obvious answer --- \emph{plug in the Pi and see} --- will cost you the
project. Consider the loop it commits you to:

\begin{center}
\begin{tikzpicture}[node distance=5mm, font=\footnotesize\sffamily,
   b/.style={draw=brewergrey, rounded corners=2pt, inner sep=4pt, minimum height=7mm}]
  \node[b] (e) {edit on laptop};
  \node[b, right=of e] (c) {copy to Pi};
  \node[b, right=of c] (m) {compile (slow)};
  \node[b, right=of m] (r) {run};
  \node[b, right=of r] (g) {squint, guess};
  \foreach \x/\y in {e/c, c/m, m/r, r/g}{
    \draw[-{Stealth[length=4pt]}, brewergrey] (\x) -- (\y);}
  \draw[-{Stealth[length=4pt]}, brewergrey] (g.south) -- ++(0,-4.5mm) -| (e.south);
  \node[font=\scriptsize\itshape, color=brewerred] at ($(m)!0.5!(r)+(0,-9mm)$)
       {two minutes per attempt};
\end{tikzpicture}
\end{center}

And every failure has three possible causes --- your logic, your wiring, the
chip --- with no way to tell them apart. You want a loop that takes 200
milliseconds and has exactly \emph{one} possible cause: your logic.

\section{The idea: program against a description, not a thing}

Your driver does not actually need ``a Raspberry Pi SPI bus''. Look at what it
truly needs: \emph{something that takes bytes out and hands bytes back}. That is
the entire requirement. A real Pi satisfies it. So does a lookup table you write
in ten lines.

So rather than have the driver depend on a concrete Pi4J class, write that
\emph{requirement} down as a type of its own --- an \textbf{interface}
\index{interface} --- and have the driver depend on the requirement.

\begin{analogy}{The wall socket}
A wall socket is an interface. Your kettle does not depend on ``the Hinkley
Point C nuclear reactor''; it depends on ``230\,V at 50\,Hz through two holes
shaped like this''. That is why the same kettle works on a generator, a solar
inverter, or a test rig in the factory.

The socket standard is a \emph{contract}: it says nothing about how the
electricity was made, only what shows up at the holes. Every appliance ever
built works with every power source ever built, because both sides code against
the socket rather than against each other.

An interface is a socket. \code{SpiLink} is the socket shape, \code{Pi4jSpiLink}
is the power station, and \code{FakeSpiLink} is the bench supply you use to test
the kettle without burning coal.
\end{analogy}

\begin{center}
\begin{tikzpicture}[font=\footnotesize\sffamily,
   b/.style={draw=brewernavy, rounded corners=2pt, inner sep=5pt, minimum width=28mm, minimum height=8mm},
   i/.style={draw=brewerblue, fill=brewerblue!8, rounded corners=2pt, inner sep=5pt, minimum width=28mm, minimum height=8mm}]
  \node[b] (drv) {\code{Rc522Driver}};
  \node[i, below=9mm of drv] (if) {\code{SpiLink}};
  \node[b, below left=12mm and 2mm of if] (fake) {\code{FakeSpiLink}};
  \node[b, below right=12mm and 2mm of if] (real) {\code{Pi4jSpiLink}};
  \draw[-{Stealth[length=4pt]}, brewernavy] (drv) --
        node[right, font=\scriptsize\itshape, color=brewergrey] {depends on} (if);
  \draw[-{Stealth[length=4pt]}, dashed, brewergrey] (fake) -- (if);
  \draw[-{Stealth[length=4pt]}, dashed, brewergrey] (real) -- (if);
  \node[font=\scriptsize\itshape, color=brewergrey, below=1mm of fake] {tests};
  \node[font=\scriptsize\itshape, color=brewergrey, below=1mm of real] {the kiosk};
\end{tikzpicture}
\end{center}

In C you would get halfway there with a struct of function pointers, and the
idea is genuinely the same. Java's version is checked by the compiler: a class
that claims to satisfy the contract but forgets a method does not compile.

\section{Build it}

\subsection{\texttt{SpiException} --- the failure channel}

SPI can fail: the device file may not exist, a transfer can error. In C you
would return $-1$ and hope the caller checks.

\begin{jcode}
package com.midnightbrewer.hardware;

/**
 * Something went wrong talking to the SPI bus.
 *
 * Deliberately a CHECKED exception (extends Exception, not RuntimeException):
 * the compiler forces every caller to either handle it or declare that it
 * throws too. Hardware failure is not an edge case you get to forget about,
 * so we make the compiler nag.
 */
public class SpiException extends Exception {

    public SpiException(String message) {
        super(message);
    }

    public SpiException(String message, Throwable cause) {
        super(message, cause);   // keeps the original error's stack trace
    }
}
\end{jcode}

\begin{concept}{Checked exceptions are a design statement}
Java has two flavours of exception. \emph{Unchecked} ones
(\code{RuntimeException}) can be ignored silently. \emph{Checked} ones must be
acknowledged --- a caller either catches it or adds \code{throws SpiException} to
its own signature.

That distinction lets you encode a rule: \emph{the bus can fail, and you may not
forget that}. The cost is that \code{throws SpiException} propagates up your call
stack, which is a little noisy. But that noise is an accurate map of everything
in your program that can fail because hardware exists, and that map turns out to
be useful.
\end{concept}

The two-argument constructor matters more than it looks. \code{super(message,
cause)} preserves the underlying error, so when Pi4J throws some
\code{IOException} deep inside, you wrap it with your own context while keeping
the original stack trace. Debug output becomes ``my error, \emph{caused by}
their error'' rather than one or the other.

\subsection{\texttt{SpiLink} --- the contract}

\begin{jcode}
package com.midnightbrewer.hardware;

/**
 * The seam between this program and an SPI bus.
 *
 * An interface is a PROMISE, not code: any class that says "implements
 * SpiLink" must supply these methods, and the compiler enforces it.
 *
 * Everything above this line -- the driver, the reader, the whole card stack
 * -- depends only on this promise. Never on Pi4J, never on a device file.
 *
 * Extends AutoCloseable so callers can use try-with-resources and the RF
 * field is never left energised by a crash.
 */
public interface SpiLink extends AutoCloseable {

    /**
     * Full-duplex SPI transfer: shift out txBuffer, shift in rxBuffer.
     *
     * SPI is inherently full duplex -- every clock pulse sends a bit AND
     * receives a bit. So this is ONE method, not a send() and a receive().
     * The API mirrors the hardware rather than pretending it is a stream.
     */
    void transfer(byte[] txBuffer, byte[] rxBuffer, int length) throws SpiException;

    /** Release the bus. Overridden to narrow the exception type. */
    @Override
    void close() throws SpiException;
}
\end{jcode}

Notice there is \textbf{no code here at all} --- no fields, no bodies. An
interface is pure vocabulary, which is exactly why it makes a good seam: there
is nothing in it that could accidentally depend on a Pi.

\begin{deeper}{The signature is a design decision}
\code{transfer(tx, rx, length)} mirrors SPI's full-duplex reality: one call,
both directions. Had it been \code{send()} and \code{receive()}, every caller
would need to know that a ``read'' is really \emph{write the address, then clock
out a dummy byte to shift the answer in} --- a hardware detail leaking into every
call site.

Design the interface around what the hardware genuinely \emph{is}, and the code
above it stays honest. Design it around what feels convenient today, and you
spend next year translating.
\end{deeper}

\subsection{\texttt{FakeSpiLink} --- the test double}

This lives in \code{src/test/java/}, so it never ships.

\begin{jcode}
package com.midnightbrewer.hardware;

import java.util.ArrayList;
import java.util.HashMap;
import java.util.List;
import java.util.Map;

/**
 * A fake SPI bus: answers from a canned table, records what it was sent.
 *
 * This is the payoff of the interface. The driver cannot tell this from a real
 * Pi -- it only ever saw the SpiLink promise, and this keeps that promise.
 */
class FakeSpiLink implements SpiLink {

    /** Register address -> value the fake chip should report. */
    private final Map<Integer, Integer> registers = new HashMap<>();

    /** Every tx buffer we were handed, so a test can assert on it. */
    final List<byte[]> sent = new ArrayList<>();

    private boolean closed = false;

    void setRegister(int address, int value) { registers.put(address, value); }

    boolean isClosed() { return closed; }

    @Override
    public void transfer(byte[] txBuffer, byte[] rxBuffer, int length) {
        sent.add(txBuffer.clone());

        // TODO 1: decode the address byte the driver sent.
        //   A READ is encoded as ((addr << 1) & 0x7E) | 0x80, so:
        //     address = (txBuffer[0] & 0x7E) >> 1
        //     isRead  = (txBuffer[0] & 0x80) != 0

        // TODO 2: if it was a read, put the canned value in rxBuffer[1].
        //   The answer arrives on the SECOND byte: the first byte out is the
        //   address, and while that goes out the chip has nothing to say yet.
        //   Default to 0 for an address nobody configured.

        // TODO 3: if it was a write, there is nothing to answer.
    }

    @Override
    public void close() { closed = true; }
}
\end{jcode}

\begin{trap}{A fake that returns nothing}
If you leave \code{transfer()} empty, \code{rxBuffer} stays all zeros, your
driver dutifully returns \code{0}, and your test asserting \code{version() ==
0x92} fails --- so you go hunting for a bug in the driver that is not there.

Worse is the opposite: a test that \emph{passes} because both sides are
trivially empty.

\textbf{A fake must be a plausible liar, not a mute.} It has to behave enough
like the real thing that the code under test is genuinely exercised.
\end{trap}

\subsection{The test}

\begin{jcode}
package com.midnightbrewer.hardware;

import org.junit.jupiter.api.Test;
import static org.junit.jupiter.api.Assertions.assertEquals;
import static org.junit.jupiter.api.Assertions.assertTrue;

class FakeSpiLinkTest {

    @Test
    void readReturnsTheCannedValue() throws SpiException {
        FakeSpiLink link = new FakeSpiLink();
        link.setRegister(0x37, 0x92);          // VersionReg answers 0x92

        // Hand-encode a READ of register 0x37, exactly as the driver will.
        byte[] tx = { (byte) (((0x37 << 1) & 0x7E) | 0x80), 0x00 };
        byte[] rx = new byte[2];
        link.transfer(tx, rx, 2);

        assertEquals(0x92, rx[1] & 0xFF);      // & 0xFF -- the signed-byte thing
        assertEquals(1, link.sent.size());
    }

    @Test
    void closeIsRecorded() throws Exception {
        FakeSpiLink link = new FakeSpiLink();
        assertTrue(!link.isClosed());
        link.close();
        assertTrue(link.isClosed());
    }
}
\end{jcode}

\begin{verifybox}
\begin{shell}
mvn test -Dtest=FakeSpiLinkTest
\end{shell}
Green --- \emph{with the reader unplugged}. Pause on that: you have just tested
hardware-facing code with no hardware in the room.
\end{verifybox}

\subsection{\texttt{Pi4jSpiLink} --- the real one}

\begin{jcode}
package com.midnightbrewer.hardware;

import com.pi4j.Pi4J;
import com.pi4j.context.Context;
import com.pi4j.io.spi.Spi;
import com.pi4j.io.spi.SpiMode;

/**
 * The real SPI bus, via Pi4J and /dev/spidev0.0.
 *
 * The ONLY class in the project that knows Pi4J exists. Everything else sees
 * the SpiLink interface. Contain your dependencies in one place and swapping
 * them later is a one-file job.
 */
public class Pi4jSpiLink implements SpiLink {

    private static final int SPI_CHANNEL = 0;         // CE0
    private static final int SPI_BAUD    = 1_000_000; // 1 MHz

    private final Context pi4j;
    private final Spi spi;

    public Pi4jSpiLink() {
        this.pi4j = Pi4J.newAutoContext();
        this.spi = pi4j.create(Spi.newConfigBuilder(pi4j)
                .id("rc522")
                .address(SPI_CHANNEL)
                .baud(SPI_BAUD)
                .mode(SpiMode.MODE_0)   // RC522 is mode 0: CPOL=0, CPHA=0
                .build());
    }

    @Override
    public void transfer(byte[] txBuffer, byte[] rxBuffer, int length)
            throws SpiException {
        try {
            spi.transfer(txBuffer, 0, rxBuffer, 0, length);
        } catch (RuntimeException e) {
            throw new SpiException("SPI transfer failed", e);
        }
    }

    @Override
    public void close() throws SpiException {
        try {
            spi.close();
            pi4j.shutdown();
        } catch (RuntimeException e) {
            throw new SpiException("failed to close SPI", e);
        }
    }
}
\end{jcode}

\paragraph{Why 1\,MHz?} The MFRC522 handles 10\,MHz on a clean PCB. On loose
jumper wires it does not. 1\,MHz is fast enough that SPI is never your bottleneck
--- the card's RF timing dominates by orders of magnitude --- and slow enough to
tolerate imperfect wiring. Optimise the thing that is actually slow, not the
thing that is easy to measure.

\paragraph{Why mode 0?} CPOL${}=0$, CPHA${}=0$, from the datasheet: clock idles
low, data sampled on the rising edge. Get this wrong and you read
plausible-looking garbage, which is the worst kind of wrong.

\section{The concept, now that you have used it}
\index{dependency injection}

\textbf{Dependency injection} is the name for what you just did:
\code{Rc522Driver} will receive its \code{SpiLink} from outside, through its
constructor, rather than constructing one itself.

\begin{jcode}
// If the driver did this, it would be welded to the Pi forever:
class Rc522Driver {
    private final SpiLink link = new Pi4jSpiLink();   // untestable
}

// Handing it in from outside means the CALLER decides:
class Rc522Driver {
    private final SpiLink link;
    Rc522Driver(SpiLink link) { this.link = link; }   // testable
}
\end{jcode}

Same code, two power sources. The kettle does not own the power station.

\begin{concept}{A rule worth internalising}
\textbf{A class should not construct the things it depends on. It should be
handed them.}

``Who decides?'' moves up to whoever assembles the program, and testability,
flexibility and clear ownership all follow from that one move. When you meet a
class that is hard to test, look first for a \code{new} in its constructor.
\end{concept}

\section{Traps}

\begin{traptable}
\code{version()} returns $-110$ & Missing \code{\& 0xFF}. \hex{0x92} as a signed byte is $-110$. \\
Fake test passes, hardware returns 0 & The fake answers unconditionally; make it decode the address. \\
\code{NoClassDefFoundError: com/sun/jna/...} & JNA missing from the classpath. Run \code{mvn dependency:copy-dependencies}. \\
Works once, fails on the second run & You did not \code{close()}; the device file stays locked. Use try-with-resources. \\
\code{Permission denied} on \code{/dev/spidev0.0} & Not in the \code{spi} group. \code{sudo usermod -aG spi,gpio \$USER}. \\
\end{traptable}

\section{Commit}

\begin{shell}
git add -A
git commit -m "M1: SpiLink interface + fake and Pi4J implementations"
\end{shell}
